\section{Постановка задачи}

Необходимо разработать экспертную систему для выбора параметров шрифтовой гарнитуры при подготовке бейджей к печати на офсетном принтере. Система должна работать в консультационном режиме: задавать пользователю вопросы о параметрах макета и условиях печати, после чего выдавать набор рекомендаций по выбору шрифта.

В рамках курсовой работы требуется:
\begin{enumerate}
    \item Описать предметную область и выделить основные критерии, влияющие на выбор гарнитуры.
    \item Сформулировать критерии выбора в обобщённом виде, без привязки к единичным частным фактам.
    \item Построить бинарное дерево решений, в котором внутренние узлы соответствуют вопросам, а листья содержат итоговые рекомендации.
    \item Реализовать полученную модель в среде CLIPS с использованием продукционных правил.
    \item Продемонстрировать работу системы на нескольких пользовательских сценариях.
\end{enumerate}

Результатом работы системы должен быть не конкретный файл шрифта, а набор типографических рекомендаций, достаточный для выбора подходящей гарнитуры и её параметров при подготовке бейджа к печати.
